\documentclass[journal]{IEEEtran}

\usepackage{amsmath}
\usepackage{amssymb}
\usepackage{graphicx}
\usepackage{cite}
\usepackage{array}
\usepackage{booktabs}
\usepackage{multirow}
\usepackage{url}
\usepackage[dvipsnames]{xcolor}

% =====================================================================
% Drafting helpers. To remove all placeholder marks in the final version,
% redefine them as no-ops:
%   \renewcommand{\needdata}[1]{}
%   \renewcommand{\needfig}[1]{}
%
%   \needdata{...} : a number that still has to be measured
%   \needfig{...}  : a figure slot; replace with \includegraphics later
%
% NOTE: keep this file pure ASCII. pdfLaTeX with IEEEtran cannot typeset
% Hangul; adding non-ASCII text here causes "Unicode character not set up"
% errors.
% =====================================================================
\newcommand{\needdata}[1]{{\color{red}\textbf{[DATA NEEDED: #1]}}}
\newcommand{\needfig}[1]{%
  \fbox{\parbox{0.45\textwidth}{\vspace{1.2cm}\centering
  {\color{red}\textbf{[FIGURE NEEDED]}}\\[4pt] #1 \vspace{1.2cm}}}}

\begin{document}

\title{FPGA-Based Time-to-Digital Converter with Dual-Phase Coarse
Timestamping and Stimulus-Independent Code-Density Calibration}

\author{Donghyun~Sung, Junki~Lee, and~Kyungtae~Lee}

\maketitle

% =====================================================================
\begin{abstract}
This paper presents a carry-chain time-to-digital converter (TDC)
implemented on a Xilinx Zynq-7000 FPGA for FMCW radar applications, together
with a systematic characterization and calibration methodology.
The delay line consists of 80 cascaded CARRY4 primitives providing 320 taps,
of which 294 codes are reachable within one 5\,ns clock period, yielding an
average bin width of 17.0\,ps.
A dual-phase coarse counter driven by $0^{\circ}$ and $180^{\circ}$ clocks
resolves coarse-timestamp ambiguity near the sampling edge.
We quantify two structural sources of nonlinearity that dominate the raw
transfer function: an intra-CARRY4 period-four bin-width pattern spanning a
$3.72\times$ range, and an entry transient at the chain input caused by
driving the first carry stage through the general-routing \texttt{CYINIT} pin
rather than the dedicated carry input.
Calibration is performed by a code-density test, and we compare two
independent stimuli---MMCM dynamic phase shifting (deterministic) and an
on-chip ring oscillator (statistical)---under a strict
calibration/validation split in which the correction table is derived from
one acquisition and evaluated on a statistically independent one.
With the ring-oscillator stimulus, peak-to-peak DNL improves from
4.03\,LSB to 0.11\,LSB and INL from 8.55\,LSB to 0.17\,LSB at unchanged
resolution.
Finally, we formulate and test a \emph{shape-invariance} hypothesis, namely
that the relative bin-width pattern is temperature invariant while only the
overall delay scale drifts, which---if valid---reduces online thermal
compensation to a single scalar correction.
\end{abstract}

\begin{IEEEkeywords}
FPGA, time-to-digital converter, carry chain, code density test,
dynamic phase shift, ring oscillator, calibration, temperature drift.
\end{IEEEkeywords}

% =====================================================================
\section{Introduction}

Time-to-digital converters (TDCs) are essential in FMCW radar, LiDAR,
positron emission tomography, and time-of-flight imaging. Modern FPGAs expose
dedicated carry-chain resources that permit tapped-delay-line (TDL) TDCs with
picosecond resolution at low development cost~\cite{song,wu,charbon}.

Practical FPGA TDL-TDCs are limited by (i) metastability when a hit arrives
near the sampling edge, (ii) strong bin-width nonuniformity arising from the
carry-chain layout, and (iii) process--voltage--temperature (PVT) drift of the
cell delay.

\subsection{Related Work and Positioning of This Work}

The individual techniques used here are established, and we state their
provenance explicitly.
Dual coarse counters clocked on opposite edges, with the fine code selecting
which counter to use, are a known solution to coarse-timestamp
metastability in TDC ASICs~\cite{atlasmdt}.
Code-density (statistical) calibration driven by an uncorrelated on-chip ring
oscillator is likewise standard~\cite{calibreview,wang19ps}.
MMCM dynamic phase shifting has been used to generate picosecond-scale
reference intervals for characterizing carry-chain TDCs~\cite{sensors2026}.
Online compensation of temperature and voltage drift by monitoring the
frequency of an on-chip ring oscillator and rescaling the delay-line
calibration was demonstrated by Bourdeauducq~\cite{bourdeauducq}.

Accordingly, this work does not claim these mechanisms as new. Our
contribution is a measurement-driven characterization and a methodological
framework built on them:

\begin{enumerate}
\item \textbf{Quantitative structural characterization} of a Zynq-7000
carry-chain delay line, separating two distinct nonlinearity mechanisms:
an intra-CARRY4 period-four pattern ($3.72\times$ bin-width range) and an
input \emph{entry transient} localized to the first three CARRY4 blocks and
attributable to \texttt{CYINIT} general-routing injection.
\item \textbf{Stimulus-independent calibration methodology.} Two physically
different stimuli (MMCM phase sweep and ring oscillator) are compared under a
calibration/validation split, so that the reported post-calibration
nonlinearity is a generalization error rather than a self-consistency
artifact.
\item \textbf{Shape-invariance hypothesis for online thermal compensation.}
Prior online schemes rescale the calibration by a single
ring-oscillator-derived factor~\cite{bourdeauducq}, but to our knowledge the
underlying premise---that the \emph{relative} bin-width pattern is temperature
invariant---has not been verified experimentally. We define the hypothesis and
a test for it.
\item An FPGA implementation and evaluation targeted at FMCW radar.
\end{enumerate}

% =====================================================================
\section{TDC Architecture}

\subsection{Overall Architecture}

Fig.~\ref{fig:overall} shows the system. A hit propagates along a carry chain
whose state is sampled by the system clock; the number of propagated taps
gives the fine time, while a coarse counter provides the clock-cycle index.

\begin{figure}[htbp]
\centering
\needfig{Overall block diagram.\\
Source: export \texttt{diagram/tdc\_system\_block.drawio} to PDF.}
\caption{Overall architecture of the proposed TDC.}
\label{fig:overall}
\end{figure}

\subsection{Carry-Chain Delay Line}

The delay line uses 80 cascaded CARRY4 primitives placed on a single slice
column by explicit \texttt{LOC}/\texttt{BEL} constraints, giving
\begin{equation}
N_{\mathrm{tap}} = 80 \times 4 = 320 .
\end{equation}
The chain is deliberately made longer than one sampling period so that the
full $T_{\mathrm{clk}}=5$\,ns is covered with margin; consequently only a
subset of taps is reachable (Section~\ref{sec:char}).

\subsection{Fine-Time Extraction}

The fine code is obtained by a population count over all sampled taps,
\begin{equation}
C_{\mathrm{fine}} = \sum_{i=0}^{319} \mathrm{Tap}_i ,
\qquad \mathrm{Tap}_i \in \{0,1\},
\end{equation}
which is inherently immune to bubbles in the thermometer code and avoids a
priority encoder. The population count is pipelined as a three-level adder
tree to close timing at 200\,MHz.

Because a measurement is only initiated when $\mathrm{Tap}_0=1$, the code
$C_{\mathrm{fine}}=0$ is structurally unreachable; the usable code range
therefore begins at unity.

\subsection{Timestamp Reconstruction}

The absolute timestamp is
\begin{equation}
T_{\mathrm{abs}} = N_{\mathrm{clk}}\,T_{\mathrm{clk}} - T_{\mathrm{fine}},
\qquad
T_{\mathrm{fine}} = \mathrm{LUT}\!\left(C_{\mathrm{fine}}\right),
\end{equation}
where the fine time is subtracted because a larger fine code corresponds to an
earlier arrival relative to the sampling edge. The LUT is a 320-entry,
13-bit on-chip ROM applied in the datapath, so calibrated timestamps are
produced at full rate without post-processing.

% =====================================================================
\section{Dual-Phase Coarse Timestamping}

\subsection{Problem}

When a hit arrives close to the sampling edge, the coarse counter may be
captured mid-transition, producing an error of one full clock period, which is
far larger than the fine-time resolution.

\begin{figure}[htbp]
\centering
\needfig{Timing diagram: metastable region near the clock edge.}
\caption{Coarse-timestamp ambiguity near a clock boundary.}
\label{fig:meta}
\end{figure}

\subsection{Implementation}

Following the dual-counter principle of~\cite{atlasmdt}, two counters are
maintained: one clocked on the rising ($0^{\circ}$) edge and one on the
falling ($180^{\circ}$) edge, the latter resynchronized into the primary
domain. The fine code selects which counter is trustworthy, since the fine
code directly indicates the position of the hit within the clock period.

A danger zone is declared when
\begin{equation}
C_{\mathrm{fine}} < T_L \quad\text{or}\quad C_{\mathrm{fine}} > T_H ,
\end{equation}
with $T_L=40$ and $T_H=220$ obtained from the measured code distribution, and
the coarse value is selected as
\begin{equation}
N_{\mathrm{clk}}=
\begin{cases}
\mathrm{Cnt}_{180}+1, & \text{danger zone},\\[2pt]
\mathrm{Cnt}_{0}, & \text{otherwise}.
\end{cases}
\end{equation}
Because the two observation points are separated by $T_{\mathrm{clk}}/2$, at
least one counter is always sampled far from its transition.

\begin{figure}[htbp]
\centering
\needfig{Dual-phase coarse timestamping concept:
$0^{\circ}$/$180^{\circ}$ counters and fine-code-based selection.}
\caption{Dual-phase coarse timestamping.}
\label{fig:dual}
\end{figure}

% =====================================================================
\section{Delay-Line Characterization}
\label{sec:char}

This section reports the measured structure of the raw transfer function,
which motivates the calibration strategy. All results are obtained by a
code-density test (Section~\ref{sec:cal}).

\subsection{Usable Code Range}

Two boundaries limit the usable range. The code $C_{\mathrm{fine}}=0$ cannot
occur because the snapshot condition requires $\mathrm{Tap}_0=1$. At the upper
end, the 320-tap chain spans approximately 5.4\,ns, exceeding the 5\,ns clock
period, so the highest taps are never reached. The measured usable range is
therefore
\begin{equation}
1 \le C_{\mathrm{fine}} \le 294 ,
\end{equation}
i.e.\ 294 codes, giving an average bin width
\begin{equation}
W_{\mathrm{LSB}} = \frac{T_{\mathrm{clk}}}{294} = 17.0~\mathrm{ps}.
\end{equation}

\subsection{Intra-CARRY4 Period-Four Pattern}

Grouping bins by their position within a CARRY4 block, $i \bmod 4$, reveals a
systematic pattern: the four outputs of a CARRY4 are not equivalent because
each is tapped at a different point of the internal carry cascade.
Table~\ref{tab:period4} lists the measured group means.

\begin{table}[htbp]
\caption{Measured Bin Width by Position within a CARRY4 Block}
\label{tab:period4}
\centering
\begin{tabular}{lcc}
\toprule
Position & Mean bin width & Relative to $W_{\mathrm{LSB}}$ \\
\midrule
$i \bmod 4 = 0$ & 10.33\,ps & 0.61 \\
$i \bmod 4 = 1$ & 31.33\,ps & 1.84 \\
$i \bmod 4 = 2$ & 17.97\,ps & 1.05 \\
$i \bmod 4 = 3$ & \phantom{0}8.43\,ps & 0.49 \\
\midrule
\multicolumn{2}{l}{Widest/narrowest ratio} & $3.72\times$ \\
\bottomrule
\end{tabular}
\end{table}

\begin{figure}[htbp]
\centering
\needfig{Scatter of bin count grouped by $i \bmod 4$ (left) and along the
chain (right).\\ Source: CARRY4 position analysis in \texttt{Histogram.py}.}
\caption{Intra-CARRY4 bin-width distribution.}
\label{fig:period4}
\end{figure}

This pattern is invariant along the chain, confirming that it originates from
the fixed internal structure of the CARRY4 primitive rather than from
placement.

\subsection{Entry Transient at the Chain Input}

The first carry stage receives the hit through the general-routing
\texttt{CYINIT} pin, whereas every subsequent stage receives its carry through
the dedicated \texttt{CIN} path. The resulting difference in edge quality
produces a transient localized at the chain input: the measured time consumed
by the first CARRY4 blocks is $1.89\times$, $1.51\times$ and $1.16\times$ the
nominal $4W_{\mathrm{LSB}}$, recovering to unity from the fourth block onward.
This transient is the single largest contributor to peak DNL.

A control experiment confirms that the mechanism is edge quality rather than
propagation delay: removing a parasitic load from the hit net reduced its
routing delay from 4227\,ps to 2143\,ps, yet peak-to-peak DNL changed only
from 4.23\,LSB to 4.20\,LSB. The constant delay offset is therefore
irrelevant, while the intrinsic \texttt{CYINIT} entry penalty remains.

\begin{figure}[htbp]
\centering
\needfig{Per-CARRY4 consumed time versus block index, showing the decay
$1.89 \rightarrow 1.51 \rightarrow 1.16 \rightarrow 1.00$.}
\caption{Entry transient at the delay-line input.}
\label{fig:entry}
\end{figure}

\subsection{Clock-Region Boundary}

Covering 5\,ns at 200\,MHz requires approximately 74 slice rows, which exceeds
the 50-row height of a 7-series clock region; the chain therefore necessarily
crosses a clock-region boundary, in our placement between taps 196 and 200.
Contrary to expectation, no significant discontinuity was observed at this
boundary: after removing the period-four component, the normalized deviation
at the boundary bins was comparable to that of interior bins elsewhere in the
chain. We conclude that, for this architecture, clock-region crossing is not a
dominant nonlinearity mechanism.

% =====================================================================
\section{Code-Density Calibration}
\label{sec:cal}

\subsection{Principle}

Under a stimulus uniformly distributed over the clock period, the number of
hits accumulated in bin $i$ is proportional to its width:
\begin{equation}
W_i = \frac{H_i}{\sum_k H_k}\, T_{\mathrm{clk}} .
\end{equation}
The calibration table stores the cumulative arrival time at the centre of each
bin,
\begin{equation}
\mathrm{LUT}[i] = \frac{T_{\mathrm{clk}}}{H_{\mathrm{tot}}}
\left( \sum_{k<i} H_k + \frac{H_i}{2} \right),
\label{eq:lut}
\end{equation}
which is monotonic by construction and is quantized to 13 bits for the on-chip
ROM.

\subsection{Calibration/Validation Split}

Evaluating a calibration on the same acquisition used to derive it is
self-consistent by construction. We therefore acquire two statistically
independent histograms, use the first to build \eqref{eq:lut}, and evaluate the
resulting linearity on the second. Denoting by $H^{\mathrm{val}}$ the
validation histogram and by $E^{\mathrm{cal}}$ the bin edges implied by the
calibration histogram, the corrected distribution is obtained by redistributing
each validation bin over the uniform grid,
\begin{equation}
M_k = \sum_i H_i^{\mathrm{val}}\,
\frac{\left| \left[E^{\mathrm{cal}}_i,E^{\mathrm{cal}}_{i+1}\right)
\cap B_k \right|}
{E^{\mathrm{cal}}_{i+1}-E^{\mathrm{cal}}_{i}} ,
\end{equation}
where $B_k$ is the $k$-th uniform time bin. The residual nonlinearity of $M_k$
is thus a generalization error.

\subsection{Stimulus 1: MMCM Dynamic Phase Shift}

A hit synchronous to a fixed clock is sampled by a phase-shifted clock; the
MMCM phase is advanced in 280 steps of $T_{\mathrm{VCO}}/56 = 17.857$\,ps,
covering exactly one 5\,ns period. This stimulus is deterministic and provides
an absolute, crystal-referenced time ruler.

\subsection{Stimulus 2: Ring Oscillator}

A 31-stage ring oscillator free-running on the FPGA fabric generates hits that
are asynchronous to the sampling clock and therefore uniformly distributed in
phase. This stimulus requires no clock manipulation and can accumulate very
large statistics in the background.

\subsection{Measured Comparison}

Table~\ref{tab:calcmp} summarizes the measured results. Both stimuli were
evaluated with the calibration/validation split described above.

\begin{table}[htbp]
\caption{Measured Calibration Performance (Validation Set)}
\label{tab:calcmp}
\centering
\begin{tabular}{lcc}
\toprule
 & MMCM phase sweep & Ring oscillator \\
\midrule
Accumulated hits (cal) & $5.6\times10^{7}$ & $1.21\times10^{8}$ \\
Accumulated hits (val) & $5.6\times10^{7}$ & $1.50\times10^{8}$ \\
Usable codes & 293 & 294 \\
$W_{\mathrm{LSB}}$ & 17.07\,ps & 17.01\,ps \\
\midrule
DNL$_{\mathrm{pp}}$ before & 4.19\,LSB & 4.03\,LSB \\
DNL$_{\mathrm{pp}}$ after  & 0.50\,LSB & \textbf{0.11\,LSB} \\
INL$_{\mathrm{pp}}$ before & 8.94\,LSB & 8.55\,LSB \\
INL$_{\mathrm{pp}}$ after  & 0.64\,LSB & \textbf{0.17\,LSB} \\
\midrule
Acquisition time & minutes (sweep) & 10--20\,min (background) \\
Absolute reference & yes (crystal) & no \\
\bottomrule
\end{tabular}
\end{table}

Both stimuli recover the same underlying bin-width structure, as expected for
two uniform samplings of the same physical delay line. The lower residual of
the ring-oscillator case follows from its larger accumulated statistics rather
than from any intrinsic advantage: the residual is dominated by the
statistical difference between the calibration and validation histograms.
The MMCM stimulus retains a distinct advantage in that its step size is
crystal-referenced and therefore provides absolute time traceability.

\begin{figure}[htbp]
\centering
\needfig{DNL before/after calibration (ring oscillator).\\
Source: output of \texttt{DNL\_INL\_codedensity.py}.}
\caption{Measured DNL before and after calibration.}
\label{fig:dnl}
\end{figure}

\begin{figure}[htbp]
\centering
\needfig{INL before/after calibration (ring oscillator).}
\caption{Measured INL before and after calibration.}
\label{fig:inl}
\end{figure}

\begin{figure}[htbp]
\centering
\needfig{Cross-stimulus validation: LUT derived from the MMCM sweep applied to
ring-oscillator data, and vice versa.}
\caption{Cross-stimulus calibration transfer.}
\label{fig:cross}
\end{figure}

\needdata{Cross-applied DNL/INL, i.e.\ MMCM-derived LUT evaluated on
ring-oscillator data and the reverse. Agreement demonstrates that the
calibration is independent of the stimulus used to obtain it.}

% =====================================================================
\section{Online Thermal Compensation and the Shape-Invariance Hypothesis}

\subsection{Motivation}

Carry-chain cell delay depends on process, supply voltage and temperature, so
a table acquired once becomes stale. Recomputing the full 294-entry table
online is expensive in both time and logic.

\subsection{Hypothesis}

Existing online schemes rescale the calibration by a single factor derived
from an on-chip ring oscillator~\cite{bourdeauducq}. This is valid only if the
\emph{relative} bin-width pattern is preserved under temperature change. We
state this explicitly:

\begin{quote}
\textbf{Shape-invariance hypothesis.} The normalized bin-width pattern
$\hat{W}_i = W_i / \overline{W}$ is invariant with temperature, while the mean
bin width $\overline{W}$ scales with the fabric delay.
\end{quote}

The rationale is that the period-four pattern and the entry transient are
consequences of the fixed internal geometry of the CARRY4 primitive and of the
routing topology, both of which are temperature independent, whereas the
absolute cell delay is not.

If the hypothesis holds, online compensation reduces to
\begin{equation}
\mathrm{LUT}_{T} = K(T)\cdot \mathrm{LUT}_{\mathrm{ref}},
\qquad
K(T) = \frac{f_{\mathrm{RO,ref}}}{f_{\mathrm{RO}}(T)} ,
\label{eq:scale}
\end{equation}
i.e.\ a single multiplication, since the ring oscillator is built from the same
fabric primitives as the delay line and therefore tracks its delay directly.
The oscillator frequency is measured against the MMCM-derived clock, which is
crystal referenced and thus temperature stable, so $K(T)$ is obtained without
any external sensor.

\subsection{Test Procedure}

The hypothesis is tested by acquiring code-density histograms at several
junction temperatures, normalizing each by its own mean, and comparing the
resulting shapes. Invariance is quantified by the residual
\begin{equation}
\varepsilon(T) = \max_i \left| \hat{W}_i(T) - \hat{W}_i(T_0) \right| ,
\end{equation}
reported in LSB, together with the fitted drift of $\overline{W}$ in
ppm/$^{\circ}$C.

\begin{figure}[htbp]
\centering
\needfig{Normalized bin-width pattern $\hat{W}_i$ at several temperatures,
overlaid. Overlap supports shape invariance.}
\caption{Test of the shape-invariance hypothesis.}
\label{fig:shape}
\end{figure}

\needdata{Code-density histograms at several junction temperatures
(e.g.\ 25, 45, 65, 85$^{\circ}$C); the shape residual $\varepsilon(T)$ in LSB;
and the drift of the mean bin width $\overline{W}$ in ppm/$^{\circ}$C.}

\subsection{Implementation}

\needdata{Resource cost and update rate of the online correction hardware
(ring-oscillator frequency counter plus scale multiplier).}

% =====================================================================
\section{Experimental Results}

\subsection{Platform}

\begin{table}[htbp]
\caption{Implementation Platform}
\centering
\begin{tabular}{lc}
\toprule
Parameter & Value \\
\midrule
FPGA device & Xilinx Zynq-7000 (Zybo) \\
Sampling clock & 200\,MHz ($T_{\mathrm{clk}}=5$\,ns) \\
Delay line & 80 $\times$ CARRY4 = 320 taps \\
Usable codes & 294 \\
Average bin width & 17.0\,ps \\
Coarse counter & dual phase ($0^{\circ}/180^{\circ}$) \\
Calibration ROM & 320 $\times$ 13\,bit \\
\bottomrule
\end{tabular}
\end{table}

\needdata{Resource utilization: LUT / FF / BRAM / DSP counts and occupancy
from the Vivado utilization report.}

\subsection{Linearity}

The measured linearity is given in Table~\ref{tab:calcmp} and
Figs.~\ref{fig:dnl}--\ref{fig:inl}. After calibration, peak-to-peak DNL and INL
are 0.11\,LSB and 0.17\,LSB respectively, i.e.\ 1.9\,ps and 2.9\,ps, at an
unchanged average bin width of 17.0\,ps.

\subsection{Dual-Phase Coarse Timestamping}

\needdata{Single-phase versus dual-phase comparison: rate of one-clock-period
timestamp discontinuities and coarse-count error rate. Suggested method:
measure a fixed time interval repeatedly and count outliers of magnitude
$T_{\mathrm{clk}}$.}

\begin{figure}[htbp]
\centering
\needfig{Timestamp stability: single-phase versus dual-phase coarse counter.}
\caption{Effect of dual-phase coarse timestamping.}
\label{fig:metares}
\end{figure}

\subsection{Single-Shot Precision}

\needdata{Standard deviation (RMS resolution) of $N$ repeated measurements of a
fixed time interval, before and after calibration.}

\subsection{Temperature Stability}

\needdata{Measurement error from 25 to 85$^{\circ}$C for three configurations:
(a) no calibration, (b) static calibration only, (c) the proposed online scale
correction.}

\begin{figure}[htbp]
\centering
\needfig{Timing error versus temperature for the three configurations.}
\caption{Temperature stability.}
\label{fig:temp}
\end{figure}

\subsection{Comparison with Prior Work}

\needdata{Comparison table against prior work: device, LSB, DNL/INL, RMS
resolution and thermal compensation method, including
\cite{bourdeauducq,wang19ps,sensors2026,calibreview}.}

% =====================================================================
\section{Conclusion}

We presented a 200\,MHz carry-chain TDC on a Zynq-7000 FPGA with dual-phase
coarse timestamping and a code-density calibration methodology evaluated under
a calibration/validation split. Measurement identifies two distinct structural
nonlinearity mechanisms---an intra-CARRY4 period-four pattern spanning
$3.72\times$ and an entry transient confined to the first three carry blocks
and caused by \texttt{CYINIT} injection---while showing that clock-region
crossing is not significant in this architecture. Calibration reduces
peak-to-peak DNL from 4.03\,LSB to 0.11\,LSB and INL from 8.55\,LSB to
0.17\,LSB at an unchanged 17.0\,ps average bin width, and the two independent
stimuli recover the same structure. Finally, we formulated the
shape-invariance hypothesis that underlies single-scalar online thermal
compensation and defined an experiment to test it.
\needdata{Concluding sentence summarizing the temperature-experiment outcome.}

% =====================================================================
\bibliographystyle{IEEEtran}
\begin{thebibliography}{99}

\bibitem{song}
J. Song, Q. An and S. Liu, ``A high-resolution time-to-digital converter
implemented in field-programmable-gate-arrays,'' \emph{IEEE Trans. Nucl. Sci.},
vol.~53, no.~1, pp.~236--241, 2006.

\bibitem{wu}
J. Wu and Z. Shi, ``The FPGA-based wave union TDC: Increasing time resolution
through multiple measurements,'' \emph{IEEE Trans. Nucl. Sci.}, vol.~55,
no.~3, pp.~1503--1507, 2008.

\bibitem{charbon}
C. Favi and E. Charbon, ``A 17 ps resolution time-to-digital converter
implemented in FPGA technology,'' in \emph{Proc. ACM/SIGDA FPGA Symp.}, 2009.

\bibitem{bourdeauducq}
S. Bourdeauducq, ``A 26 ps RMS time-to-digital converter core for Spartan-6
FPGAs,'' arXiv:1303.6840, 2013.

\bibitem{calibreview}
``Calibration methods for time-to-digital converters,'' \emph{Sensors}, 2023.
\needdata{verify authors, volume, issue and pages}

\bibitem{wang19ps}
``A 19 ps precision and 170 MSamples/s time-to-digital converter implemented in
FPGA with online calibration,'' \emph{Applied Sciences}, vol.~12, no.~7,
p.~3649, 2022.

\bibitem{sensors2026}
``Application and comparison of FPGA-based carry chain TDC and DDMTD schemes in
high-precision time synchronization,'' \emph{Sensors}, vol.~26, no.~3, 2026.
\needdata{verify authors and pages}

\bibitem{atlasmdt}
``Development of a time-to-digital converter ASIC for the upgrade of the ATLAS
monitored drift tube detector,'' arXiv:1708.03692, 2017.

\bibitem{xilinx_clock}
Xilinx Inc., ``7 series FPGAs clocking resources user guide (UG472).''

\bibitem{xilinx_clb}
Xilinx Inc., ``7 series FPGAs configurable logic block user guide (UG474).''

\end{thebibliography}

\end{document}
